\section{Conclusion}
\label{sec:conclusion}

We have formulated the portability of a comparative scheduler ranking,
$R(W, M, L)$, as an object of simulation-based scheduler selection and
comparison, distinct from both workload characterization and
single-scheduler evaluation, and designed, froze, and executed a paired,
multi-source trace-driven simulation protocol -- a mechanism-diverse fixed
13-policy scheduler panel, a calibrated six-region operating grid, an
explicit metric contract, and an uncertainty/reversal-detection methodology
-- against $9{,}360$ structurally validated configurations before
interpreting any comparative outcome. Across three measurably different
workload sources (\S\ref{sec:workloads}), comparative rankings are portable
in aggregate (Kendall's tau-b never below $0.55$) but source-pair-dependent:
two sources agree almost perfectly, while the third is consistently less
portable against either at every operating region (\S\ref{sec:results}).
Supported practical reversals are uncommon ($3.6\%$ of pairwise
comparisons) and concentrated: all involve the same policy pair and
BurstGPT, at or above the pre-knee region (\S\ref{sec:reversals}).
Exact-order recovery of the full-corpus ranking requires most or all of the
40-window budget for two of three sources, whereas identifying the single
best policy requires far fewer windows for those same two sources
(\S\ref{sec:sample-complexity}).

Two post-campaign extensions show that portability is substantially weaker
across evaluation metrics than across sources (median $\tau_b = 0.419$;
\S\ref{sec:results-cross-metric}) and only partly robust to the SLO
definition (\S\ref{sec:results-slo}). RQ3 remains a pilot-scale engineering
validation, not a transfer conclusion (\S\ref{sec:results-rq3-pilot}). A
selected physical vLLM/GPU validation (RQ6) did not reproduce the largest
simulated reversal, while a paired stable control retained its qualitative
ordering -- a selected-case simulator-fidelity boundary, not a general
estimate of hardware reversal prevalence (\S\ref{sec:real-system}).

Ranking portability, in this simulation study, is therefore neither
universal nor uniformly fragile: it is real, but conditional on the workload
source, operating region, evaluation metric, SLO definition, and model
fidelity context under which a comparison was established. The appropriate
lesson is not to abandon simulation, but to report simulation-derived
scheduler-selection claims with their experimental conditions and validity
boundaries explicit. LSSP is released as a public, independently
reproducible artifact (\S\ref{sec:artifact}) so that each result can be
verified and extended.
